\appendix
\section*{Appendix A — Mathematical Foundations}
\addcontentsline{toc}{section}{Appendix A — Mathematical Foundations}

\section{Preliminaries}

\subsection{State Manifold}
Let $\mathcal{M}$ denote the Fun–Sphere narrative manifold. Each point
$p \in \mathcal{M}$ carries scalar fields $f(p)$, vector fields $v(p)$,
curvature annotations $\kappa(p)$, and spectral fields $\lambda(p)$.
These fields collectively define the geometric state structure used
throughout the engine.

\subsection{Operator Group}
Operators are reversible transformations
\[
O : \mathcal{S} \rightarrow \mathcal{S},
\qquad
O^{-1} \text{ exists},
\]
acting on the state space $\mathcal{S}$. Operators form a group under
composition, with identity $I$ and inverses $O^{-1}$.

\subsection{Holonomy and Curvature}
Holonomy arises from non–commutative operator pairs. For operators
$O_i, O_j$,
\[
\Delta_{\mathrm{hol}}(O_i, O_j)
    = (O_j \circ O_i) - (O_i \circ O_j),
\]
defines the fundamental curvature tensor of the engine. Curvature
accumulates along operator paths and induces frame shifts.

\subsection{Laplacian Geometry}
The discrete Laplacian is defined as
\[
L = D - W,
\]
where $D$ is the degree matrix and $W$ is the curvature–aware adjacency
matrix. The Laplacian is symmetric and positive semi–definite, admitting
eigenpairs $(\lambda_k, u_k)$.

\subsection{Spectral Ecology}
Eigenvectors $u_k$ represent monster families, while eigenvalues
$\lambda_k$ represent ecological energy levels. Curvature perturbs the
Laplacian, producing spectral drift:
\[
\lambda_k' = \lambda_k + \delta \lambda_k,
\qquad
u_k' = u_k + \delta u_k.
\]

\subsection{Reversible Quest Algebra}
A reversible quest is an operator path
\[
\gamma = (O_1, O_2, \ldots, O_n),
\]
with forward traversal $\gamma(s)$ and reverse traversal
$\gamma^{-1}(s)$. State reversibility holds, but curvature and frame
shifts do not reverse.

\section{Definitions}

\begin{definition}[State Space]
The state space $\mathcal{S}$ is the set of all valid Fun–Sphere engine
states, each containing scalar, vector, spectral, and curvature fields.
\end{definition}

\begin{definition}[Reversible Operator]
An operator $O$ is reversible if there exists $O^{-1}$ such that
\[
O^{-1}(O(s)) = s
\]
for all $s \in \mathcal{S}$.
\end{definition}

\begin{definition}[Holonomy Tensor]
The holonomy tensor $H$ is defined by
\[
H_{ij} = \Delta_{\mathrm{hol}}(O_i, O_j),
\]
and measures curvature induced by operator non–commutativity.
\end{definition}

\begin{definition}[Quest Curvature]
For a path $\gamma = (O_1, \ldots, O_n)$,
\[
K(\gamma) = \sum_{i < j} H_{ij}
\]
is the total curvature accumulated along the quest.
\end{definition}

\section{Theorems}

\begin{theorem}[Positive Semi–Definiteness of the Laplacian]
For any vector $x$,
\[
x^\top L x \ge 0.
\]
\end{theorem}

\begin{proof}
Since $L = D - W$ with $W$ symmetric and $D$ diagonal with non–negative
entries, the quadratic form expands to
\[
x^\top L x = \sum_{i < j} w_{ij} (x_i - x_j)^2 \ge 0.
\]
\end{proof}

\begin{theorem}[State Reversibility of Quests]
For any quest path $\gamma$ and state $s$,
\[
\gamma^{-1}(\gamma(s)) = s.
\]
\end{theorem}

\begin{proof}
By definition,
\[
\gamma(s) = O_n(\cdots O_2(O_1(s))\cdots),
\]
and
\[
\gamma^{-1}(s) = O_1^{-1}(O_2^{-1}(\cdots O_n^{-1}(s)\cdots)).
\]
Since $O_i^{-1}(O_i(s)) = s$ for all operators, the composition of
inverses in reverse order restores the original state.
\end{proof}

\begin{theorem}[Non–Reversibility of Curvature]
Let $\gamma$ be a quest path. Then
\[
K(\gamma^{-1}) \neq -K(\gamma),
\]
and in general curvature does not reverse.
\end{theorem}

\begin{proof}
Curvature depends on operator non–commutativity:
\[
H_{ij} = (O_j \circ O_i) - (O_i \circ O_j).
\]
Reversing the path changes the order of operators but does not negate
the holonomy tensor. Thus curvature accumulates independently of state
reversibility.
\end{proof}

\section{Discussion}
This appendix establishes the formal mathematical foundations required
for the operator algebra, holonomy curvature, Laplacian geometry,
spectral ecology, and reversible quest dynamics presented in the main
preprint. Subsequent appendices provide detailed proofs, extended
derivations, and formal verification artifacts.
